\section{Section 7: The main theorem}

Throughout this section, for complex numbers $u_0, u_1$ with
$|u_0| = |u_1| = 1$, we define
$R(u_0, u_1) = \{(u_0, u_1),\; (1, u_0^* u_1)\}$.

\begin{lemma}
  \label{lem:s7_1}
  \lean{TwoControl.section7_lemma_7_1}
  \leanok
  Suppose $u_0, u_1$ are complex numbers with $|u_0| = |u_1| = 1$.
  For a $2$-qubit unitary $W$ and a $3$-qubit unitary $U$,
  if $CC(\operatorname{Diag}(u_0, u_1)) = W^{AB} \cdot U$,
  then there exists a $2$-qubit unitary $V$ such that
  $CC(\operatorname{Diag}(1, u_0^* u_1)) = V^{AB} \cdot U$.
\end{lemma}

\begin{proof}
  Set $V = C(\operatorname{Diag}(1, u_0^*)) \cdot W$.
  Then $V$ is a $2$-qubit unitary.
  Expand $V^{AB} \cdot U$ using the hypothesis
  $W^{AB} \cdot U = CC(\operatorname{Diag}(u_0, u_1))$ to obtain
  $V^{AB} \cdot U
    = C(\operatorname{Diag}(1, u_0^*))^{AB}
      \cdot CC(\operatorname{Diag}(u_0, u_1))$.
  Both factors are diagonal in the computational basis,
  so the product is verified entry-by-entry:
  the diagonal entries combine to give
  $CC(\operatorname{Diag}(1, u_0^* u_1))$.
\end{proof}

\begin{lemma}
  \label{lem:s7_2}
  \lean{TwoControl.section7_lemma_7_2}
  \leanok
  Suppose $u_0, u_1$ are complex numbers with $|u_0| = |u_1| = 1$.
  If there exists a product of at most four elements of
  $G_2 \setminus G_1$ and any number of elements of $G_1$ that equals
  $CC(\operatorname{Diag}(u_0, u_1))$, then either:
  \begin{enumerate}
    \item there exist $2$-qubit unitaries $U_1, U_2, U_3, U_4$
      and $(u_2, u_3) \in R(u_0, u_1)$ such that
      $U_1^{BC}\, U_2^{AC}\, U_3^{AB}\, U_4^{BC}
        = CC(\operatorname{Diag}(u_2, u_3))$,
      or
    \item there exist $2$-qubit unitaries $U_1, U_2, U_3, U_4$
      and $(u_2, u_3) \in R(u_0, u_1)$ such that
      $U_1^{AC}\, U_2^{BC}\, U_3^{AC}\, U_4^{BC}
        = CC(\operatorname{Diag}(u_2, u_3))$.
  \end{enumerate}
\end{lemma}

\begin{proof}
  \uses{lem:s7_1}
  The proof proceeds in four steps.

  (\emph{Step 1.})
  Each element of $G_1$ can be absorbed into an adjacent
  element of $G_2 \setminus G_1$, reducing to at most four
  factors from $G_2 \setminus G_1$.

  (\emph{Step 2.})
  By padding with identity factors if fewer than four
  $G_2 \setminus G_1$ factors remain, write the product
  as exactly four canonical factors $V_1 V_2 V_3 V_4$,
  where each $V_i$ is of the form $U^{XY}$ for some pair $XY$.

  (\emph{Step 3.})
  Using the swap conjugation identities (Lemma~A.12), reduce
  to the case where $V_4$ has label $BC$.
  This may change the diagonal parameters from $(u_0, u_1)$
  to $(1, u_0^* u_1)$ via Lemma~\ref{lem:s7_1},
  so the new parameters $(u_2, u_3) \in R(u_0, u_1)$.

  (\emph{Step 4.})
  Using Lemma~A.13 and the swap conjugation identities,
  case-analyse over the labels of $V_1, V_2, V_3$ to reduce
  to one of the two canonical forms stated in the conclusion.
\end{proof}

\begin{lemma}
  \label{lem:s7_3}
  \lean{TwoControl.section7_lemma_7_3}
  \leanok
  Suppose $u_0, u_1$ are complex numbers with $|u_0| = |u_1| = 1$.
  Suppose $(u_2, u_3) \in R(u_0, u_1)$.
  If $u_2 = u_3$ or $u_2 u_3 = 1$,
  then $u_0 = u_1$ or $u_0 u_1 = 1$.
\end{lemma}

\begin{proof}
  Since $(u_2, u_3) \in R(u_0, u_1)$, either
  $(u_2, u_3) = (u_0, u_1)$ or $(u_2, u_3) = (1, u_0^* u_1)$.
  Cross with the hypothesis $u_2 = u_3$ or $u_2 u_3 = 1$
  to obtain four subcases.
  In the case $(u_2, u_3) = (u_0, u_1)$, the conclusion holds
  immediately.
  In the case $(u_2, u_3) = (1, u_0^* u_1)$ with $u_2 = u_3$,
  we get $u_0^* u_1 = 1$, so $u_0 u_1 = |u_0|^2 u_0 u_1 / u_0 u_1$;
  using $|u_0| = 1$ gives $u_0 = u_1$.
  In the case $(u_2, u_3) = (1, u_0^* u_1)$ with $u_2 u_3 = 1$,
  we get $u_0^* u_1 = 1$, hence $u_0 = u_1$.
\end{proof}

\begin{theorem}
  \label{thm:s7_4}
  \lean{TwoControl.section7_theorem_7_4}
  \leanok
  Suppose $u_0, u_1$ are complex numbers with $|u_0| = |u_1| = 1$.
  There exists a product of at most four elements of
  $G_2 \setminus G_1$ and any number of elements of $G_1$ that equals
  $CC(\operatorname{Diag}(u_0, u_1))$ if and only if
  either $u_0 = u_1$ or $u_0 u_1 = 1$.
\end{theorem}

\begin{proof}
  \uses{lem:s7_2, lem:s5_1, lem:s6_4, lem:s7_3}
  (\emph{Forward direction.})
  By Lemma~\ref{lem:s7_2}, the hypothesis reduces to one of two
  canonical forms with parameters $(u_2, u_3) \in R(u_0, u_1)$.
  In form~(1), $U_1^{BC}\, U_2^{AC}\, U_3^{AB}\, U_4^{BC}
    = CC(\operatorname{Diag}(u_2, u_3))$,
  so by Lemma~\ref{lem:s5_1} either $u_2 = u_3$ or $u_2 u_3 = 1$.
  In form~(2), $U_1^{AC}\, U_2^{BC}\, U_3^{AC}\, U_4^{BC}
    = CC(\operatorname{Diag}(u_2, u_3))$,
  so by Lemma~\ref{lem:s6_4} either $u_2 = u_3$ or $u_2 u_3 = 1$.
  In either case, Lemma~\ref{lem:s7_3} lifts the conclusion to
  $u_0 = u_1$ or $u_0 u_1 = 1$.

  (\emph{Backward direction.})
  If $u_0 = u_1$ or $u_0 u_1 = 1$, Lemma~\ref{lem:s5_1} (backward)
  provides a decomposition
  $U_1^{BC}\, U_2^{AC}\, U_3^{AB}\, U_4^{BC}
    = CC(\operatorname{Diag}(u_0, u_1))$,
  each factor being an element of $G_2$.
\end{proof}

\begin{corollary}
  \label{cor:s7_5}
  \lean{TwoControl.section7_corollary_7_5}
  \leanok
  For a $1$-qubit unitary $U$, there exists a product of
  at most four elements of $G_2 \setminus G_1$ and any number
  of elements of $G_1$ that equals $CC(U)$ if and only if
  either the eigenvalues of $U$ are equal or $\det(U) = 1$.
\end{corollary}

\begin{proof}
  \uses{thm:s7_4}
  By the spectral theorem (Theorem~A.3), write
  $U = V \operatorname{Diag}(u_0, u_1) V^\dagger$
  for a unitary $V$ and eigenvalues $u_0, u_1$ with $|u_0| = |u_1| = 1$.
  Then $CC(U) = (I \otimes I \otimes V)\,
    CC(\operatorname{Diag}(u_0, u_1))\,
    (I \otimes I \otimes V^\dagger)$.
  Since $I \otimes I \otimes V \in G_1$, conjugating by it
  does not change the number of $G_2 \setminus G_1$ factors needed.
  Apply Theorem~\ref{thm:s7_4} to $CC(\operatorname{Diag}(u_0, u_1))$
  to conclude: the decomposition exists if and only if
  $u_0 = u_1$ or $u_0 u_1 = 1$.
  The condition $u_0 = u_1$ is equivalent to the eigenvalues being
  equal, and $u_0 u_1 = \det(U) = 1$.
\end{proof}